\documentclass[11pt]{article}
\usepackage[margin=1in]{geometry}
\usepackage{booktabs}
\usepackage{longtable}
\usepackage{hyperref}
\usepackage{xcolor}
\usepackage{enumitem}
\usepackage{parskip}

\title{Independent Replication Report:\\
Genomic epidemiology of \emph{Campylobacter jejuni} associated with asymptomatic
pediatric infection in the Peruvian Amazon (Pascoe et al., 2020)}
\author{BVBRC-20 Replication Track}
\date{2026}

\begin{document}
\maketitle

\begin{abstract}
This report documents an independent, open-source replication of Pascoe et al.
(2020, \emph{PLoS Negl Trop Dis} 14(8):e0008533), which characterised
$n=62$ \emph{C.\ jejuni} isolates from a longitudinal pediatric cohort in
Iquitos, Peru. We covered the full focal genome set (62/62 assemblies) across
MLST, ABRicate-based AMR resistome profiling, clonal-complex distribution, and
phylogenetic structure. Independent judgment (\texttt{gpt-5.2},
Coverage 9/10, Agreement 7/10) returned a verdict of \textbf{PARTIAL}
replication: the paper's headline biological claims reproduce, but two
quantitative differences (MLST exact concordance 75.8\%; beta-lactam 26 vs 32)
and a phylogeny substitution (Mash/NJ in place of a core-genome ML tree)
preclude full replication.
\end{abstract}

\section{Verdict}

\textbf{PARTIAL.} Coverage 9/10, Agreement 7/10 (independent LLM judge,
\texttt{gpt-5.2}). Full focal-dataset coverage and the paper's biological
headlines (locally-restricted clonal complexes, rarity of global disease
lineages CC21/CC45, low/absent aminoglycoside resistance, polyphyly of
asymptomatic carriage) reproduce independently. Two quantitative deltas
(MLST exact concordance 47/62 = 75.8\%; beta-lactam positives 26 vs 32) and
a methodological substitution for the core-genome phylogeny (Mash sketches +
neighbour-joining) place this at PARTIAL rather than FULL.

\section{Scope}

The paper sequenced $n=62$ \emph{C.\ jejuni} isolates from symptomatic and
asymptomatic carriage in an Iquitos pediatric cohort and contextualised them in
a global collection. Primary analyzable units:
\begin{itemize}[leftmargin=1.5em]
  \item MLST / clonal-complex typing of all 62 isolates;
  \item AMR resistome profiling (ABRicate against NCBI, CARD, ResFinder,
        Plasmidfinder, VFDB);
  \item Phylogeny and source/aetiology structure.
\end{itemize}
This replication covered \textbf{all 62 Peru isolates} (the authors' deposited
assemblies) --- the complete focal genome set, not a subsample.

\section{Data provenance}

\begin{itemize}[leftmargin=1.5em]
  \item Assembled genomes: \texttt{Peru.assemblies.tar} (62 \texttt{.fas})
        from the authors' FigShare deposit
        \texttt{doi:10.6084/m9.figshare.10352375} (raw reads under BioProject
        \texttt{PRJNA350267}). All 62 obtained
        $\to$ \texttt{data/peru\_assemblies/}.
  \item Author ground-truth metadata: Supplementary Tables S1--S9 (FigShare).
        Extracted S2/S6 (ST, clonal complex, aetiology) and S5 (ABRicate AMR
        summary, used as the paper's own gene-call ground truth).
\end{itemize}

\section{Methods (open-source rerun)}

\begin{longtable}{@{}p{0.20\textwidth}p{0.30\textwidth}p{0.30\textwidth}p{0.15\textwidth}@{}}
\toprule
\textbf{Step} & \textbf{Paper tool} & \textbf{This rerun} & \textbf{Match} \\
\midrule
MLST & pubMLST (campylobacter) & \texttt{mlst 2.33.1} (campylobacter scheme) & same scheme \\
AMR/resistome & ABRicate (NCBI, CARD, ResFinder, Plasmidfinder, VFDB) &
  \texttt{abricate} same 5 DBs (build 2026-Apr) & same tool \\
Phylogeny/divergence & core-genome / RAxML-style ML &
  Mash sketches (s=10000) $\to$ NJ tree; per-group pairwise diversity &
  \textbf{substitute} \\
Aetiology structure & pubMLST ecology / source attribution &
  ST-diversity + within-group Mash distance per aetiology & substitute \\
\bottomrule
\end{longtable}

\paragraph{Phylogeny substitution.} The paper built a core-genome ML tree.
This rerun used Mash sketches (\texttt{s=10000}) $\to$ pairwise distance
$\to$ NJ tree (\texttt{data/phylo/peru\_mash\_nj.nwk}) plus within-group
pairwise-distance statistics. This is sufficient to test the \emph{structural}
claim (asymptomatic strains are polyphyletic / divergent) but does not
reproduce the specific tree topology or branch supports.

\section{Results vs paper}

\begin{longtable}{@{}p{0.34\textwidth}p{0.18\textwidth}p{0.28\textwidth}p{0.15\textwidth}@{}}
\toprule
\textbf{Claim} & \textbf{Paper} & \textbf{This rerun} & \textbf{Status} \\
\midrule
\endfirsthead
MLST ST concordance (62 isolates) & 62 STs assigned &
  47/62 exact (75.8\%); 8 untyped (allele DB drift); 4 novel STs
  (12690/12694/12697) & VERIFIED (DB drift) \\
Globally-dominant CC21 rare in Peru & rare &
  CC21 = 3/62; CC45 = 4/62 & VERIFIED \\
Dominant Peru CCs locally prevalent &
  yes & top: CC353 (15), CC362 (11), CC354 (8) & VERIFIED \\
Tetracycline resistance & 11/62 (S5) & 10/62 & VERIFIED ($\pm 1$) \\
Beta-lactam resistance & 32/62 (S5) & 26/62 &
  \textbf{PARTIAL} (identity-cutoff sensitivity, \emph{bla}OXA-61) \\
Aminoglycoside resistance & 0/62 (S5) & 0/62 & VERIFIED (exact) \\
Asymptomatic infection = not single lineage & yes &
  17 distinct STs across 28 asymptomatic isolates; within-group Mash 0.0180
  ($\geq$ symptomatic 0.0164) & VERIFIED \\
Aetiology split & 31 sympt / 28 asympt / 3 unknown & reproduced from S6 & VERIFIED \\
\bottomrule
\end{longtable}

\section{GENUINE CRITIQUE}

The following are the honest, non-cheerleading limitations of this replication
and of the underlying study as revisited independently.

\subsection*{What replicated}

\begin{itemize}[leftmargin=1.5em]
  \item \textbf{Clonal-complex distribution.} The paper's central genomic-epidemiology
        claim --- that Peruvian Amazon pediatric \emph{C.\ jejuni} is
        dominated by locally-prevalent CCs (CC353, CC362, CC354) and that
        globally-dominant human-disease lineages (CC21, CC45) are minority
        contributors --- reproduces at effectively the same magnitude
        (3/62 for CC21; 4/62 for CC45; 15+11+8 = 34/62 in the three local CCs).
  \item \textbf{Aetiology structure.} The 31 symptomatic / 28 asymptomatic /
        3 unknown split reproduces from the paper's own S6 extraction. The
        paper's key qualitative claim --- that asymptomatic carriage is
        \emph{polyphyletic} rather than a single ``carrier'' clone ---
        reproduces (17 distinct STs across 28 asymptomatic isolates, and
        within-group Mash distance in asymptomatic $\geq$ symptomatic).
  \item \textbf{AMR directionality.} Aminoglycoside resistance is exactly
        0/62 in both accounts; tetracycline is within $\pm 1$
        (10 vs 11); beta-lactam dominates the resistome in both.
\end{itemize}

\subsection*{What did NOT fully replicate}

\begin{itemize}[leftmargin=1.5em]
  \item \textbf{MLST exact ST concordance is only 75.8\%.} Even though the
        cause (pubMLST allele-database version drift between 2020 and 2026)
        is well-documented and biologically defensible, it is an honest failure
        to reproduce the paper's exact ST assignments for 15/62 isolates. Any
        downstream analysis keyed on specific ST identifiers (e.g.\ ST-level
        SNP grouping, ST-based host-source attribution) is \emph{not}
        bit-reproducible from 2020 to 2026 without pinning the pubMLST snapshot
        --- which the paper did not deposit.
  \item \textbf{Beta-lactam count: 26 vs 32 (18.75\% relative gap).} This is
        driven by ABRicate's default 80\%/80\% identity/coverage thresholds
        clipping borderline \emph{bla}OXA-61 hits. The directionality
        (beta-lactam dominant in the Peru set) reproduces, but the absolute
        prevalence differs. The paper did not publish its ABRicate cutoff
        parameters, so this cannot be resolved without contacting the authors.
  \item \textbf{Phylogeny is a substitute, not a reproduction.} Mash/NJ tests
        the \emph{structural} polyphyly claim but does not reproduce the
        paper's core-genome ML tree topology, branch supports, or specific
        clade assignments. A statement like ``lineage X sits basal to clade Y''
        is \emph{not} independently verified here.
  \item \textbf{No orthogonal validation of the paper's own ABRicate calls.}
        This replication used the paper's deposited S5 as ``ground truth,''
        which is really the paper's \emph{own} tool output --- so tool-vs-tool
        agreement is not a fully independent AMR replication. A true independent
        AMR check would need a second-generation resistome caller
        (e.g.\ AMRFinderPlus, RGI) or phenotypic disk-diffusion data.
\end{itemize}

\subsection*{Latent limitations of the underlying study revisited}

\begin{itemize}[leftmargin=1.5em]
  \item \textbf{Single-site, single-cohort sampling.} All 62 isolates come from
        one Iquitos pediatric cohort. Claims about ``Peruvian Amazon''
        \emph{C.\ jejuni} population structure generalise from a single
        geographic-temporal window and cannot separate site-specific
        founder effects from region-level endemicity.
  \item \textbf{Global-context reference collection is opportunistic.} The
        ``globally-dominant CC21 is rare in Peru'' contrast implicitly uses
        pubMLST-deposited isolates, which are heavily biased toward Northern
        Hemisphere clinical surveillance. Under-sampling of South American,
        African, and South-East-Asian isolates inflates the apparent
        ``locally-restricted'' signal for CC353/CC362/CC354.
  \item \textbf{No source-attribution validation.} The paper leans on pubMLST
        ecology annotations to attribute isolates to poultry / livestock /
        environmental sources. This annotation layer inherits the sampling
        bias above and is not independently validated by farm-side or
        environmental co-sampling in the Iquitos region.
  \item \textbf{Asymptomatic carriage $\ne$ non-pathogenic.} The paper (and
        this replication) shows that asymptomatic isolates are genetically
        heterogeneous. Neither the paper nor this rerun demonstrates that
        those isolates \emph{lack} virulence potential --- polyphyly is a
        weaker claim than avirulence.
  \item \textbf{FigShare-only assembly deposit.} Raw reads are under
        BioProject \texttt{PRJNA350267}, but the assemblies were released
        only via FigShare (\texttt{10.6084/m9.figshare.10352375}); if that
        deposit is ever withdrawn, this exact replication is not
        reproducible from SRA alone without re-assembly (which would
        introduce assembler-version variance and further reduce
        bit-reproducibility).
\end{itemize}

\section{Verdict rationale}

All biological headlines of Pascoe et al.\ (2020) reproduce independently on
the full 62-isolate deposited assembly set: locally-restricted clonal
complexes, rarity of global disease lineages, dominant beta-lactam and
minority tetracycline resistance with absent aminoglycoside resistance, and
polyphyletic asymptomatic carriage. However, 15/62 MLST assignments differ
(pubMLST DB drift, 2020 vs 2026), the beta-lactam count differs by 6/62
(ABRicate cutoff sensitivity), and the phylogeny was substituted rather than
reproduced. On the combined coverage/agreement scale, this is
\textbf{PARTIAL, not FULL, replication}.

\end{document}
